\documentclass[12pt]{article}
\usepackage[utf8]{inputenc}
\usepackage[T1]{fontenc}
\usepackage[margin=1in]{geometry}
\usepackage{setspace}
\usepackage{xcolor}
\usepackage[colorlinks=true, urlcolor=blue!60!black, linkcolor=blue!60!black]{hyperref}
\setlength{\parindent}{2em}
\setlength{\parskip}{1.25em}
\renewcommand{\baselinestretch}{1.0}
\begin{document}
\begin{center}
{\LARGE \textit{Bivens v.\ Six Unknown Named Agents of Federal Bureau of Narcotics}, 403 U.S.\ 388 (1971)\par}
\vspace{0.8em}
{\large Case Brief --- Four-Page Edition\par}
\end{center}
\vspace{1em}

\section{Memory Jogger}

Webster Bivens alleged that federal narcotics agents entered his apartment without a warrant or probable cause, arrested him in front of his family, and strip-searched him. A 6--3 Court held that he could seek damages directly under the Fourth Amendment, finding no special factors counseling hesitation and no congressional direction to use another remedy (\href{https://supreme.justia.com/cases/federal/us/403/388/}{Justia}).

\section{Detailed Case Facts}

On November 26, 1965, agents of the Federal Bureau of Narcotics, acting under claim of federal authority, entered Webster Bivens's Brooklyn apartment without a warrant or probable cause and arrested him for alleged narcotics violations. According to his complaint, they manacled him in front of his wife and children, threatened to arrest the whole family, and searched the apartment ``from stem to stern.'' They took him to the federal courthouse in Brooklyn, where he was interrogated, booked, and subjected to a visual strip search.

\section{Procedural History}

In 1967 Bivens sued in the Eastern District of New York for the humiliation, embarrassment, and mental suffering caused by the warrantless entry, search, and arrest with unreasonable force, seeking \$15,000 from each agent. Because 42 U.S.C.\ \S~1983 reaches only state officials, no statute authorized a damages suit against federal officers. The district court dismissed for failure to state a federal cause of action and on immunity grounds; the Second Circuit affirmed on the ground that the Fourth Amendment creates no damages remedy. The Supreme Court reversed and remanded without reaching immunity.

\section{Judicial Votes}

Reversed and remanded, \textbf{6--3} (decided June 21, 1971).
\begin{itemize}
\item \textbf{Majority:} Justice Brennan, joined by Justices Douglas, Stewart, White, and Marshall.
\item \textbf{Concurring in the judgment:} Justice Harlan.
\item \textbf{Dissents:} Chief Justice Burger; Justice Black; Justice Blackmun.
\end{itemize}

\section{Holding}

A violation of the Fourth Amendment by federal agents acting under color of federal authority gives rise to a federal cause of action for damages directly under the Constitution, absent ``special factors counseling hesitation'' or an explicit congressional declaration remitting injured persons to an equally effective alternative remedy.

\section{Analysis of Opinions}

\textbf{Majority (Brennan).} Resolving the question left open in \textit{Bell v.\ Hood}, Brennan rejected the government's view that Bivens should be left to state trespass law. First, federal power ``does not disappear like a magic gift when it is wrongfully used,'' and an agent acting in the government's name is far more dangerous than a private trespasser, so the Fourth Amendment independently limits federal power. Second, state law may be inconsistent with or hostile to Fourth Amendment interests: one may bar the door against a private intruder, but an agent's claim of authority makes resistance futile or criminal, and state law cannot define the reach of federal authority. Third, damages are the ordinary remedy for invasions of liberty, and under \textit{Bell} and \textit{Marbury} courts may use traditional remedies; no special factors counseled hesitation, and Congress had provided no substitute.

\textbf{Harlan (concurring).} Since federal-question jurisdiction lets courts enjoin constitutional violations, it lets them award damages when ``necessary'' or ``appropriate.'' Injunctions cannot undo a completed search, the United States is immune, and the exclusionary rule does nothing for an innocent victim never prosecuted.

\textbf{Burger (dissenting).} Creating damages remedies is Congress's job; he attacked the exclusionary rule and urged a statutory compensation tribunal. \textbf{Black (dissenting).} Section 1983 shows Congress knows how to create such suits; dockets were already choked. \textbf{Blackmun (dissenting).} The Court engaged in ``judicial legislation'' that would bring ``another avalanche of new federal cases.''

\textbf{Later.} The Court extended \textit{Bivens} only in \textit{Davis v.\ Passman} (1979) and \textit{Carlson v.\ Green} (1980). \textit{Ziglar v.\ Abbasi} (2017) called expansion a ``disfavored judicial activity'' and formalized a two-step test asking whether a claim arises in a ``new context'' and whether ``special factors'' counsel hesitation; \textit{Egbert v.\ Boule} (2022) asked whether there is \textit{any} reason to think Congress better equipped; holding that even a less effective alternative process forecloses relief; and \textit{Goldey v.\ Fields} (2025) unanimously denied a remedy for a federal prisoner's Eighth Amendment excessive-force claim.

\section{Examples: Future Applications}

\begin{enumerate}
\item \textbf{Same-Side: A DEA Home Raid.} A warrantless, abusive raid by successor narcotics agents tracks \textit{Bivens} itself, avoiding a ``new context.''
\item \textbf{Same-Side: Denied Emergency Care.} A federal prisoner left to die of an asthma attack fits \textit{Carlson v.\ Green}, an existing context.
\item \textbf{Opposite-Side: A TSA Checkpoint.} Excessive force and retaliation by a TSA screener arise in a new context with aviation-security concerns and alternative remedies.
\item \textbf{Opposite-Side: A Demoted Federal Engineer.} Civil-service remedies foreclose a First Amendment claim, as in \textit{Bush v.\ Lucas} (1983).
\item \textbf{Fence-Sitter: An FBI Fraud Raid.} FBI agents conduct a \textit{Bivens}-like raid in a bank-fraud case; whether a different agency and statute make a ``new context'' is unsettled.
\end{enumerate}

\section{Critique}

\textbf{Marbury versus separation of powers.} \textit{Bivens} pits Marshall's maxim that every right needs a remedy against the modern view that creating damages actions is a legislative choice. \textbf{The \S~1983 asymmetry.} Because Congress never enacted a federal counterpart to \S~1983, \textit{Abbasi}, \textit{Egbert}, and \textit{Goldey} leave citizens with far weaker remedies against federal agents than against state police for identical acts, pushing plaintiffs toward the FTCA, which bars punitive damages and jury trials.

\section{Key Quotations}

\begin{quote}
``The present case involves no special factors counseling hesitation in the absence of affirmative action by Congress. ... For we have here no explicit congressional declaration that persons injured by a federal officer's violation of the Fourth Amendment may not recover money damages from the agents, but must instead be remitted to another remedy, equally effective in the view of Congress.''
\end{quote}
\noindent --- Justice Brennan, for the Court, 403 U.S.\ at 396--97.

\begin{quote}
``Finally, assuming Bivens' innocence of the crime charged, the `exclusionary rule' is simply irrelevant. For people in Bivens' shoes, it is damages or nothing.''
\end{quote}
\noindent --- Justice Harlan, concurring in the judgment, 403 U.S.\ at 410.

\begin{quote}
``An agent acting---albeit unconstitutionally---in the name of the United States possesses a far greater capacity for harm than an individual trespasser exercising no authority other than his own.''
\end{quote}
\noindent --- Justice Brennan, for the Court, 403 U.S.\ at 392.

\end{document}
